\chapter{Justificación social}
\label{cap:justificacion}

El proyecto busca beneficiar principalmente a los trabajadores de la
agencia que atienden consultas y reservas. Contar con una fuente de
información sobre los paquetes, las fechas y los cupos puede facilitar su
trabajo y ayudar a que todos sigan el mismo procedimiento antes de
confirmar una solicitud. De esta forma se pretende reducir el tiempo que se
dedica a revisar datos o solucionar diferencias entre reservas.

Los clientes también pueden resultar beneficiados. Al solicitar un viaje
esperan recibir información clara sobre la disponibilidad, el precio y los
servicios incluidos. Si la agencia maneja datos actualizados y registra
correctamente cada cambio, será más sencillo dar respuestas oportunas y
evitar confusiones al momento de confirmar una reserva.

De manera indirecta, una mejor organización de los pasajeros confirmados
puede servir para coordinar el transporte, los guías y el hospedaje. Por
tratarse de una propuesta académica, no se afirmará que el sistema
aumentará las ventas o reducirá gastos en una cantidad determinada hasta
que se realicen pruebas y exista información que permita demostrarlo.
